\documentclass[10pt,a4paper]{amsart}
\usepackage{amsmath,amssymb,amsthm,graphicx}
% For dealing with references we use the comment environment
\usepackage{verbatim}
\newenvironment{reference}{\comment}{\endcomment}
%\newenvironment{reference}{}{}
\newenvironment{slogan}{\comment}{\endcomment}
\newenvironment{history}{\comment}{\endcomment}
% For commutative diagrams we use Xy-pic
\usepackage[all]{xy}

% We use 2cell for 2-commutative diagrams.
\xyoption{2cell}
\UseAllTwocells
\theoremstyle{plain}
\newtheorem{theorem}[subsection]{Theorem}
\newtheorem{proposition}[subsection]{Proposition}
\newtheorem{lemma}[subsection]{Lemma}

\theoremstyle{definition}
\newtheorem{definition}[subsection]{Definition}
\newtheorem{example}[subsection]{Example}
\newtheorem{exercise}[subsection]{Exercise}
\newtheorem{situation}[subsection]{Situation}

\theoremstyle{remark}
\newtheorem{remark}[subsection]{Remark}
\newtheorem{remarks}[subsection]{Remarks}

\numberwithin{equation}{subsection}
% Macros
%
\def\lim{\mathop{\mathrm{lim}}\nolimits}
\def\colim{\mathop{\mathrm{colim}}\nolimits}
\def\Spec{\mathop{\mathrm{Spec}}}
\def\Hom{\mathop{\mathrm{Hom}}\nolimits}
\def\Ext{\mathop{\mathrm{Ext}}\nolimits}
\def\SheafHom{\mathop{\mathcal{H}\!\mathit{om}}\nolimits}
\def\SheafExt{\mathop{\mathcal{E}\!\mathit{xt}}\nolimits}
\def\Sch{\mathit{Sch}}
\def\Mor{\mathop{\mathrm{Mor}}\nolimits}
\def\Ob{\mathop{\mathrm{Ob}}\nolimits}
\def\Sh{\mathop{\mathit{Sh}}\nolimits}
\def\NL{\mathop{N\!L}\nolimits}
\def\CH{\mathop{\mathrm{CH}}\nolimits}
\def\proetale{{pro\text{-}\acute{e}tale}}
\def\etale{{\acute{e}tale}}
\def\QCoh{\mathit{QCoh}}
\def\Ker{\mathop{\mathrm{Ker}}}
\def\Im{\mathop{\mathrm{Im}}}
\def\Coker{\mathop{\mathrm{Coker}}}
\def\Coim{\mathop{\mathrm{Coim}}}

% Boxtimes
%
\DeclareMathSymbol{\boxtimes}{\mathbin}{AMSa}{"02}

%
% Macros for moduli stacks/spaces
%
\def\QCohstack{\mathcal{QC}\!\mathit{oh}}
\def\Cohstack{\mathcal{C}\!\mathit{oh}}
\def\Spacesstack{\mathcal{S}\!\mathit{paces}}
\def\Quotfunctor{\mathrm{Quot}}
\def\Hilbfunctor{\mathrm{Hilb}}
\def\Curvesstack{\mathcal{C}\!\mathit{urves}}
\def\Polarizedstack{\mathcal{P}\!\mathit{olarized}}
\def\Complexesstack{\mathcal{C}\!\mathit{omplexes}}
% \Pic is the operator that assigns to X its picard group, usage \Pic(X)
% \Picardstack_{X/B} denotes the Picard stack of X over B
% \Picardfunctor_{X/B} denotes the Picard functor of X over B
\def\Pic{\mathop{\mathrm{Pic}}\nolimits}
\def\Picardstack{\mathcal{P}\!\mathit{ic}}
\def\Picardfunctor{\mathrm{Pic}}
\def\Deformationcategory{\mathcal{D}\!\mathit{ef}}
\usepackage{cleveref}
\crefname{originallemma}{original lemma}{original lemmas}
\Crefname{originallemma}{Original lemma}{Original lemmas}
\makeatletter
\newlabel{lemma-quasi-compact-dense-open-connected-at-infinity-Noetherian}{{0BQF}{source}{}{sourcepointer-1}{}}
\newlabel{lemma-quasi-compact-dense-open-connected-at-infinity-Noetherian@cref}{{[originallemma][0][]0BQF}{[1][0][]source}{}{}{}}
\makeatother
\crefname{originallemma}{original lemma}{original lemmas}
\Crefname{originallemma}{Original lemma}{Original lemmas}
\makeatletter
\newlabel{local-cohomology-lemma-depth-2-connected-punctured-spectrum}{{0BLR}{source}{}{sourcepointer-2}{}}
\newlabel{local-cohomology-lemma-depth-2-connected-punctured-spectrum@cref}{{[originallemma][0][]0BLR}{[1][0][]source}{}{}{}}
\makeatother
\crefname{originallemma}{original lemma}{original lemmas}
\Crefname{originallemma}{Original lemma}{Original lemmas}
\makeatletter
\newlabel{algebra-lemma-apply-grothendieck}{{0337}{source}{}{sourcepointer-3}{}}
\newlabel{algebra-lemma-apply-grothendieck@cref}{{[originallemma][0][]0337}{[1][0][]source}{}{}{}}
\makeatother
\crefname{originallemma}{original lemma}{original lemmas}
\Crefname{originallemma}{Original lemma}{Original lemmas}
\makeatletter
\newlabel{lemma-extend-S2}{{0EY6}{source}{}{sourcepointer-4}{}}
\newlabel{lemma-extend-S2@cref}{{[originallemma][0][]0EY6}{[1][0][]source}{}{}{}}
\makeatother
\crefname{originaltheorem}{original theorem}{original theorems}
\Crefname{originaltheorem}{Original theorem}{Original theorems}
\makeatletter
\newlabel{algebra-theorem-nullstellensatz}{{00FV}{source}{}{sourcepointer-5}{}}
\newlabel{algebra-theorem-nullstellensatz@cref}{{[originaltheorem][0][]00FV}{[1][0][]source}{}{}{}}
\makeatother
\crefname{originallemma}{original lemma}{original lemmas}
\Crefname{originallemma}{Original lemma}{Original lemmas}
\makeatletter
\newlabel{algebra-lemma-finite-free-given-residue-field-extension}{{07NE}{source}{}{sourcepointer-6}{}}
\newlabel{algebra-lemma-finite-free-given-residue-field-extension@cref}{{[originallemma][0][]07NE}{[1][0][]source}{}{}{}}
\makeatother
\crefname{originallemma}{original lemma}{original lemmas}
\Crefname{originallemma}{Original lemma}{Original lemmas}
\makeatletter
\newlabel{more-morphisms-lemma-slice-smooth}{{057G}{source}{}{sourcepointer-7}{}}
\newlabel{more-morphisms-lemma-slice-smooth@cref}{{[originallemma][0][]057G}{[1][0][]source}{}{}{}}
\makeatother
\crefname{originallemma}{original lemma}{original lemmas}
\Crefname{originallemma}{Original lemma}{Original lemmas}
\makeatletter
\newlabel{algebra-lemma-section-smooth}{{05D5}{source}{}{sourcepointer-8}{}}
\newlabel{algebra-lemma-section-smooth@cref}{{[originallemma][0][]05D5}{[1][0][]source}{}{}{}}
\makeatother
\crefname{originallemma}{original lemma}{original lemmas}
\Crefname{originallemma}{Original lemma}{Original lemmas}
\makeatletter
\newlabel{algebra-lemma-completion-flat}{{00MB}{source}{}{sourcepointer-9}{}}
\newlabel{algebra-lemma-completion-flat@cref}{{[originallemma][0][]00MB}{[1][0][]source}{}{}{}}
\makeatother
\crefname{originallemma}{original lemma}{original lemmas}
\Crefname{originallemma}{Original lemma}{Original lemmas}
\makeatletter
\newlabel{lemma-purity-power-series-over-depth2}{{0EYA}{source}{}{sourcepointer-10}{}}
\newlabel{lemma-purity-power-series-over-depth2@cref}{{[originallemma][0][]0EYA}{[1][0][]source}{}{}{}}
\makeatother
\title{Lemma purity smooth over depth2 in pione.tex}
\author{The Stacks Project Authors}
\date{stacks-project:a04446e57ec1fbc252a871afcec7752fb2807b14; GNU Free Documentation License 1.2 (as published in the Stacks Project COPYING)}
\begin{document}
\maketitle
\noindent\textit{Editorial note.} All mathematical statement and proof text is reproduced verbatim from the author's chapter pione.tex of the Stacks Project, fixed at commit a04446e57ec1fbc252a871afcec7752fb2807b14. The pinned preamble is supplied as its own file; the wrapper class, the xr-hyper wiring and the \textbackslash{}externaldocument list are omitted because the assembler supplies the class.\par\medskip
\noindent\textit{Editorial note.} Numbering inside this standalone excerpt is editorial, not the full Stacks book print numbering.\par\medskip
\noindent\textit{Editorial note.} The candidate is from the extractor's stacks-candidates line, not from the earlier judged initial\_corpus line; the H2 judgement and its reason below were made in this module by reading the author's own statement and complete proof in the fixed checkout. Original candidate key: stacks:pione.tex:lemma-purity-smooth-over-depth2.\par\medskip
\noindent\textit{Editorial note.} Reference labels in the proof resolve to navigation pointers to the original author source; these pointers do not reproduce proofs of the prerequisites.\par\medskip
\noindent\textit{Allowed standard prerequisites (editorial).}\begin{itemize}
\item Prerequisites named by the author inside the counted span. The counted span is reproduced verbatim, including the author's own citations. Its prerequisites are the earlier results of the same fixed Stacks Project version that the author cites there and in the chapter's lead-in.
\end{itemize}\medskip
\ifdefined\hypertarget\hypertarget{sourcepointer-1}{}\fi
\noindent\textit{Source reference.} Original lemma 0BQF: sources/pione.tex, source lines 2019--2019. Navigation pointer used in the counted span; the span contains the pointer only and the cited result is not reproduced inside it. Registered as a non-load-bearing navigation stub in the established form for Stacks material. No bibliography citation is involved.\par\medskip
\ifdefined\hypertarget\hypertarget{sourcepointer-2}{}\fi
\noindent\textit{Source reference.} Original lemma 0BLR: sources/local-cohomology.tex, source lines 328--328. Navigation pointer used in the counted span; the span contains the pointer only and the cited result is not reproduced inside it. Registered as a non-load-bearing navigation stub in the established form for Stacks material. No bibliography citation is involved.\par\medskip
\ifdefined\hypertarget\hypertarget{sourcepointer-3}{}\fi
\noindent\textit{Source reference.} Original lemma 0337: sources/algebra.tex, source lines 46374--46374. Navigation pointer used in the counted span; the span contains the pointer only and the cited result is not reproduced inside it. Registered as a non-load-bearing navigation stub in the established form for Stacks material. No bibliography citation is involved.\par\medskip
\ifdefined\hypertarget\hypertarget{sourcepointer-4}{}\fi
\noindent\textit{Source reference.} Original lemma 0EY6: sources/pione.tex, source lines 5220--5220. Navigation pointer used in the counted span; the span contains the pointer only and the cited result is not reproduced inside it. Registered as a non-load-bearing navigation stub in the established form for Stacks material. No bibliography citation is involved.\par\medskip
\ifdefined\hypertarget\hypertarget{sourcepointer-5}{}\fi
\noindent\textit{Source reference.} Original theorem 00FV: sources/algebra.tex, source lines 6524--6524. Navigation pointer used in the counted span; the span contains the pointer only and the cited result is not reproduced inside it. Registered as a non-load-bearing navigation stub in the established form for Stacks material. No bibliography citation is involved.\par\medskip
\ifdefined\hypertarget\hypertarget{sourcepointer-6}{}\fi
\noindent\textit{Source reference.} Original lemma 07NE: sources/algebra.tex, source lines 44675--44675. Navigation pointer used in the counted span; the span contains the pointer only and the cited result is not reproduced inside it. Registered as a non-load-bearing navigation stub in the established form for Stacks material. No bibliography citation is involved.\par\medskip
\ifdefined\hypertarget\hypertarget{sourcepointer-7}{}\fi
\noindent\textit{Source reference.} Original lemma 057G: sources/more-morphisms.tex, source lines 10775--10775. Navigation pointer used in the counted span; the span contains the pointer only and the cited result is not reproduced inside it. Registered as a non-load-bearing navigation stub in the established form for Stacks material. No bibliography citation is involved.\par\medskip
\ifdefined\hypertarget\hypertarget{sourcepointer-8}{}\fi
\noindent\textit{Source reference.} Original lemma 05D5: sources/algebra.tex, source lines 38407--38407. Navigation pointer used in the counted span; the span contains the pointer only and the cited result is not reproduced inside it. Registered as a non-load-bearing navigation stub in the established form for Stacks material. No bibliography citation is involved.\par\medskip
\ifdefined\hypertarget\hypertarget{sourcepointer-9}{}\fi
\noindent\textit{Source reference.} Original lemma 00MB: sources/algebra.tex, source lines 22956--22956. Navigation pointer used in the counted span; the span contains the pointer only and the cited result is not reproduced inside it. Registered as a non-load-bearing navigation stub in the established form for Stacks material. No bibliography citation is involved.\par\medskip
\ifdefined\hypertarget\hypertarget{sourcepointer-10}{}\fi
\noindent\textit{Source reference.} Original lemma 0EYA: sources/pione.tex, source lines 5807--5807. Navigation pointer used in the counted span; the span contains the pointer only and the cited result is not reproduced inside it. Registered as a non-load-bearing navigation stub in the established form for Stacks material. No bibliography citation is involved.\par\medskip
\begin{lemma}
\label{lemma-purity-smooth-over-depth2}
Let $f : X \to S$ be a morphism of schemes. Let $U \subset X$
be an open subscheme. Assume
\begin{enumerate}
\item $f$ is smooth,
\item $S$ is Noetherian,
\item for $s \in S$ with $\text{depth}(\mathcal{O}_{S, s}) \leq 1$
we have $X_s = U_s$,
\item $U_s \subset X_s$ is dense for all $s \in S$.
\end{enumerate}
Then $\textit{F\'Et}_X \to \textit{F\'Et}_U$ is an equivalence.
\end{lemma}

\begin{proof}
The functor is fully faithful by
Lemma \ref{lemma-quasi-compact-dense-open-connected-at-infinity-Noetherian}
combined with Local Cohomology,
Lemma \ref{local-cohomology-lemma-depth-2-connected-punctured-spectrum}
(plus an application of
Algebra, Lemma \ref{algebra-lemma-apply-grothendieck}
to check the depth condition).

\medskip\noindent
Let $\pi : V \to U$ be a finite \'etale morphism. Let $Y \to X$
be the finite morphism constructed in Lemma \ref{lemma-extend-S2}.
We have to show that $Y \to X$ is finite \'etale.
To show that this is true for all points $x \in X$ mapping to a
given point $s \in S$ we may perform a base change by a flat
morphism $S' \to S$ of Noetherian schemes such that $s$ is
in the image. This follows from the compatibility of the
construction in Lemma \ref{lemma-extend-S2} with flat base change.

\medskip\noindent
After enlarging $U$ we may assume $U \subset X$ is
the maximal open over which $Y \to X$ is finite \'etale.
Let $Z \subset X$ be the complement of $U$.
To get a contradiction, assume $Z \not = \emptyset$.
Let $s \in S$ be a point in the image of $Z \to S$
such that no strict generalization of $s$ is in the image.
Then after base change to $\Spec(\mathcal{O}_{S, s})$
we see that $S = \Spec(A)$ with $(A, \mathfrak m, \kappa)$
a local Noetherian ring of depth $\geq 2$ and $Z$
contained in the closed fibre $X_s$
and nowhere dense in $X_s$. Choose a closed point $z \in Z$.
Then $\kappa(z)/\kappa$ is finite (by the Hilbert Nullstellensatz, see
Algebra, Theorem \ref{algebra-theorem-nullstellensatz}).
Choose a finite flat morphism $(S', s') \to (S, s)$ of local schemes
realizing the residue field extension $\kappa(z)/\kappa$, see
Algebra, Lemma \ref{algebra-lemma-finite-free-given-residue-field-extension}.
After doing a base change by $S' \to S$ we reduce to the case
where $\kappa(z) = \kappa$.

\medskip\noindent
By More on Morphisms, Lemma \ref{more-morphisms-lemma-slice-smooth}
there exists a locally closed subscheme $S' \subset X$ passing through $z$
such that $S' \to S$ is \'etale at $z$. After performing the base change
by $S' \to S$, we may assume there is a section $\sigma : S \to X$
such that $\sigma(s) = z$. Choose an affine neighbourhood
$\Spec(B) \subset X$ of $s$. Then $A \to B$ is a smooth ring
map which has a section $\sigma : B \to A$. Denote $I = \Ker(\sigma)$
and denote $B^\wedge$ the $I$-adic completion of $B$.
Then $B^\wedge \cong A[[x_1, \ldots, x_d]]$ for some $d \geq 0$, see
Algebra, Lemma \ref{algebra-lemma-section-smooth}.
Observe that $d > 0$ since otherwise we see that $X \to S$
is \'etale at $z$ which would imply that $z$ is a generic point of
$X_s$ and hence $z \in U$ by assumption (4).
Similarly, if $d > 0$, then $\mathfrak m B^\wedge$ maps into
$U$ via the morphism $\Spec(B^\wedge) \to X$.
It suffices prove $Y \to X$ is finite \'etale after base change
to $\Spec(B^\wedge)$. Since $B \to B^\wedge$ is flat
(Algebra, Lemma \ref{algebra-lemma-completion-flat})
this follows from Lemma \ref{lemma-purity-power-series-over-depth2}
and the uniqueness in the construction of $Y \to X$.
\end{proof}

\end{document}
